\section{Background}
\label{sec:background}

\subsection{From Stabilization to Geometry}

\CodeWM{} uses a JEPA-style latent prediction setup
\citep{assran2023ijepa}.  An online encoder maps the pre-edit code state
$s_t$ to $z_t = f_\theta(s_t)$, an action encoder maps the edit action
$a_t$ into the same latent space, and a predictor estimates the next
target state $\hat{z}_{t+1} = g_\phi(z_t, h_\psi(a_t))$.  The target is
produced by a stop-gradient target encoder
$z_{t+1}^{\mathrm{tar}} = f_{\bar{\theta}}(s_{t+1})$ whose parameters are
maintained by EMA.

The first \CodeWM{} paper found that the target update rule dominates
training stability in this small-model regime.  With fast EMA tracking,
the target encoder collapses and consecutive states become almost
identical under the target representation.  The leading indicator is:
\begin{equation}
    \copycos =
    \mathbb{E}_t\left[
    \cos\left(f_{\bar{\theta}}(s_t),
              f_{\bar{\theta}}(s_{t+1})\right)\right].
    \label{eq:copycos}
\end{equation}
When this value approaches 1, simply copying the previous target state is
nearly as good as predicting the next one.  Holding the target encoder
near-static with EMA decay $0.99999$ prevents that collapse and enables
longer training~\citep{akbulut2026collapse}.

The present work keeps that stabilized regime fixed and examines the
encoder's internal geometry.  The central distinction is between
\emph{training stability} and \emph{transition signal preservation}.
The first paper showed how to keep training alive; this paper shows that
a stable final representation can still discard much of the information
needed for explicit state-space deltas.

\subsection{Looped Encoders}

The encoder is weight-shared: a single transformer block is applied
repeatedly.  If $X_0$ is the embedded code sequence, then:
\begin{equation}
    X_\ell = \mathcal{T}_{\theta}(X_{\ell-1}),
    \quad \ell = 1,\ldots,L.
    \label{eq:looped_encoder}
\end{equation}
The original stabilized recipe used $L=6$ loops.  Weight sharing gives
depth at low parameter count, but it also means the same smoothing
operator is repeatedly applied.  Phase~8 probes indicate that this
recursion progressively compresses the edit signal.

\subsection{Prediction, Direction, Aux Losses, and VICReg}

The stabilized baseline optimizes prediction to the target embedding and
directional consistency in latent displacement:
\begin{align}
    \mathcal{L}_{\mathrm{pred}}
    &= 1 - \cos(\hat{z}_{t+1}, \operatorname{sg}[z_{t+1}^{\mathrm{tar}}]), \\
    \mathcal{L}_{\mathrm{dir}}
    &= 1 - \cos(\hat{z}_{t+1} - z_t,\,
                \operatorname{sg}[z_{t+1}^{\mathrm{tar}}] - z_t).
\end{align}
Phase~9 adds auxiliary supervision at intermediate loop outputs:
\begin{equation}
    \mathcal{L}
    =
    \mathcal{L}_{\mathrm{final}}
    + \lambda_{\mathrm{aux}}
      \sum_{\ell \in \mathcal{A}_{\mathrm{loops}}}
      \mathcal{L}_{\mathrm{aux}}^{(\ell)}.
    \label{eq:aux_loss}
\end{equation}
The first successful Phase~9 variant uses predictive auxiliary losses
rather than contrastive auxiliary losses.  In the current recipe,
$\mathcal{A}_{\mathrm{loops}}=\{1,2\}$ and
$\lambda_{\mathrm{aux}}=0.3$.

The later readout-collapse diagnostic makes the final regularizer
load-bearing.  The current best configuration replaces the old
readout regularization with a VICReg-style objective~\citep{bardes2022vicreg}
that combines invariance, per-dimension variance, and covariance
decorrelation terms.  This matters because cosine prediction and lift can
look strong even when the pooled representation occupies only a tiny
subspace.  We therefore report lift alongside effective rank and KNN
retrieval over validation embeddings.

\paragraph{Related latent-dynamics objectives.}
Concurrent work from the original JEPA
group~\citep{straightening2026} adds a curvature penalty to
JEPA-style latent planning, encouraging linearly interpolable
trajectories.  JEPA-Reasoner~\citep{jepareasoner2025}
reinterprets looped encoder blocks as latent reasoning steps, with
intermediate supervision as ``latent chain-of-thought'' training.
Both framings are complementary to the present recipe: VICReg can be
viewed as a discrete curvature penalty, and predictive auxiliary losses
at intermediate loops are a form of supervision on latent reasoning
steps.
